\section[Notebook 1971-8: $G \times G$ modulo the diagonal]{Notebook 1971-8: $G \times G$ modulo the normal closure of the diagonal\protect\verifie{Notebook1971_8}}\label{sec:quillen-1971-8}

\emph{Context.} \carnet{1971}{8} opens, on 2 June 1971, with problems left by Quillen's work on the group completion of a topological monoid $M$. He studies a space built from $M \times M$ with the diagonal action of $M$, which should be the group completion, and local coefficient systems on it.

\emph{The page.} On pages~7 and~8 he shows that for a local system compatible with the action of $M$, a loop of $M$ embedded diagonally acts trivially. The action of $\pi_1(M) \times \pi_1(M)$ on a fibre therefore kills the diagonal subgroup, and he concludes that the action is abelian, since the quotient of $G \times G$ by the normal subgroup generated by the diagonal $\Delta G$ is the abelianization $G_{\mathrm{ab}}$. He deduces that the identity component of the space has an abelian fundamental group.

\emph{Intuition.} Modulo the diagonal, $(x, y)$ is equivalent to $(xy^{-1}, 1)$, so the quotient is a quotient of $G$. The diagonal is not normal in $G \times G$: conjugating $(b, b)$ by $(a, 1)$ gives $(aba^{-1}, b)$, and dividing by $(b, b)$ gives $(aba^{-1}b^{-1}, 1)$. So the normal subgroup generated by the diagonal contains the commutators in the first factor, and the quotient of $G$ obtained is exactly $G_{\mathrm{ab}}$.

\begin{encadre}
\begin{theoreme}[Quillen, notebook 1971-8, p.~8]\label{thm:quillen-1971-8}
Let $G$ be a group and $N$ the normal subgroup of $G \times G$ generated by the diagonal $\{(g, g)\}$. The map $(x, y) \mapsto [x]\,[y]^{-1}$ induces an isomorphism from $(G \times G)/N$ onto the abelianization $G_{\mathrm{ab}} = G/[G, G]$.
\end{theoreme}
\end{encadre}

\begin{proof}
The map $\varphi : (x, y) \mapsto [x][y]^{-1}$ is a homomorphism because $G_{\mathrm{ab}}$ is abelian, it is onto since $\varphi(x, 1) = [x]$, and its kernel is a normal subgroup containing the diagonal, hence containing $N$. Conversely let $(x, y)$ be in the kernel, so that $c = xy^{-1}$ lies in $[G, G]$. The set of $c \in G$ with $(c, 1) \in N$ is a subgroup of $G$, and it contains every commutator, by the computation above: $(a, 1)(b, b)(a, 1)^{-1}(b, b)^{-1} = (aba^{-1}b^{-1}, 1)$. Hence $(c, 1) \in N$, and $(x, y) = (c, 1)(y, y) \in N$.
\end{proof}

\emph{Verification in Lean.} The file \href{\leanblobmain/Notebook1971_8.lean}{\texttt{Notebook1971\_8.\allowbreak{}lean}} defines \texttt{diagClosure G} with \texttt{Subgroup.normalClosure} and the homomorphism \texttt{toAb G} to \texttt{Abelianization G}. The computation with commutators is \texttt{commutator\_le\_comap}; the kernel and the surjectivity are \texttt{ker\_toAb} and \texttt{toAb\_surjective}. Theorem~\ref{thm:quillen-1971-8} is the isomorphism \texttt{quotientEquiv}, with its formula \texttt{quotientEquiv\_mk}. The topological statements of the page, on local systems and on the fundamental group of the space, are not formalised. The axioms used are the three usual axioms of mathlib.

\litterature{The group-theoretic statement is elementary and the page states it as known. Its consequence on the page is, in published form, part of the group completion theorem: the fundamental group of a group-like $H$-space is abelian (McDuff--Segal, \emph{Invent.\ Math.}~31 (1976), as recalled). We did not find the statement in mathlib. The edition does not present it as new.}
